\documentclass[11pt,a4paper]{article}

% ----------------- Paquetes -----------------
\usepackage[utf8]{inputenc}
\usepackage[T1]{fontenc}
\usepackage{geometry}
\geometry{left=2cm,right=2cm,top=2.2cm,bottom=2.2cm}
\usepackage[table]{xcolor}
\usepackage{tcolorbox}
\tcbuselibrary{skins,breakable}
\usepackage{enumitem}
\usepackage{tabularx}
\usepackage{booktabs}
\usepackage{titlesec}
\usepackage{fancyhdr}
\usepackage{hyperref}
\usepackage{array}

% ----------------- Colores -----------------
\definecolor{navy}{RGB}{22,50,79}
\definecolor{teal}{RGB}{25,114,120}
\definecolor{griscaja}{RGB}{245,247,249}
\definecolor{azulsuave}{RGB}{232,240,250}
\definecolor{nivelbajo}{RGB}{144,238,144}
\definecolor{nivelmedio}{RGB}{255,224,102}
\definecolor{nivelalto}{RGB}{255,173,92}
\definecolor{nivelcritico}{RGB}{255,92,96}

\hypersetup{
  colorlinks=true,
  linkcolor=teal,
  urlcolor=teal,
  pdfauthor={Joaquin Herrera Gamez},
  pdftitle={Reporte ejecutivo de riesgos - Clinica privada}
}

% ----------------- Cajas de contenido -----------------
\newtcolorbox{decisionbox}[1][Decision requerida]{%
  enhanced, breakable,
  colback=azulsuave, colframe=navy,
  fonttitle=\bfseries, title=#1
}

% ----------------- Encabezado y pie -----------------
\pagestyle{fancy}
\fancyhf{}
\fancyhead[L]{\small\color{navy} Gestion de riesgos -- Clinica privada}
\fancyhead[R]{\small\color{navy} Reporte al Directorio}
\fancyfoot[L]{\small\color{navy} Joaquin Herrera Gamez -- DNI 45360092}
\fancyfoot[R]{\small\color{navy} Pagina \thepage\ de \pageref{LastPage}}
\renewcommand{\headrulewidth}{0.4pt}
\usepackage{lastpage}

% ----------------- Formato de titulos -----------------
\titleformat{\section}{\Large\bfseries\color{navy}}{\thesection.}{0.6em}{}
\titlespacing*{\section}{0pt}{2.5ex plus 1ex minus .2ex}{1.2ex plus .2ex}
\titleformat{\subsection}{\large\bfseries\color{teal}}{\thesubsection.}{0.6em}{}

\newcommand{\metric}[3]{%
\begin{minipage}[t]{0.235\linewidth}
\colorbox{#1}{\parbox[c][1.35cm][c]{0.92\linewidth}{\centering\color{navy}\textbf{\Large #2}\\[-1mm]\scriptsize #3}}
\end{minipage}}

\begin{document}

% ---------- Portada / encabezado ----------
\begin{center}
\colorbox{navy}{\parbox{\textwidth}{\centering
\vspace{3mm}
{\color{white}\Huge\bfseries Reporte Ejecutivo de Riesgos}\\[2mm]
{\color{white}\large Clinica privada -- Registro inicial de riesgos en SimpleRisk}\\[2mm]
{\color{white}\small Dirigido al Directorio \quad$\vert$\quad Septiembre 2026 \quad$\vert$\quad Escenario ficticio (TP Seguridad de Sistemas)}
\vspace{3mm}
}}
\end{center}

\vspace{4mm}

\section*{Resumen ejecutivo}
\addcontentsline{toc}{section}{Resumen ejecutivo}

Se realiz\'o un an\'alisis de riesgos de seguridad de la informaci\'on porque no se
contemplaban formalmente los escenarios en los que los sistemas de la cl\'inica
podr\'ian verse comprometidos. Se identificaron 7 riesgos espec\'ificos del contexto
de la organizaci\'on. El m\'as urgente es R07: el personal utiliza contrase\~nas
d\'ebiles o reutilizadas, f\'acilmente vulnerables tanto por un atacante externo como
por alguien dentro de la organizaci\'on, con riesgo de acceso indebido a informaci\'on
sensible. Ya se encuentran en planificaci\'on 3 acciones concretas para reducir los
riesgos de mayor nivel.

\vspace{2mm}
\noindent
\metric{griscaja}{7}{riesgos evaluados}\hfill
\metric{nivelcritico!55}{1}{critico}\hfill
\metric{nivelalto!55}{2}{altos}\hfill
\metric{azulsuave}{3}{planes definidos}

\vspace{4mm}

\begin{decisionbox}
Autorizar el inicio de los 3 planes de acci\'on ya definidos y priorizar la
implementaci\'on de una pol\'itica de contrase\~nas obligatoria junto con
autenticaci\'on multifactor (MFA), por tratarse del riesgo de mayor nivel (R07,
Cr\'itico).
\end{decisionbox}

\vspace{4mm}

\section*{Top 5 riesgos por nivel}
\addcontentsline{toc}{section}{Top 5 riesgos por nivel}

\begin{center}
\small
\renewcommand{\arraystretch}{1.6}
\begin{tabularx}{\textwidth}{@{} >{\centering\arraybackslash}p{1.1cm} X >{\centering\arraybackslash}p{2.1cm} >{\centering\arraybackslash}p{1.9cm} @{}}
\toprule
\textbf{ID} & \textbf{Riesgo} & \textbf{Prob.\ $\times$ Imp.} & \textbf{Nivel} \\
\midrule
R07 & Contrase\~nas d\'ebiles o reutilizadas & 4 $\times$ 4 = 16 & \cellcolor{nivelcritico}\textbf{Cr\'itico} \\
R01 & Acceso indebido a historias cl\'inicas por personal sin necesidad de uso & 3 $\times$ 4 = 12 & \cellcolor{nivelalto}\textbf{Alto} \\
R06 & Equipamiento m\'edico (IoT) como puerta de entrada a la red & 3 $\times$ 4 = 12 & \cellcolor{nivelalto}\textbf{Alto} \\
R03 & Phishing dirigido al personal administrativo y m\'edico & 2 $\times$ 4 = 8 & \cellcolor{nivelmedio}\textbf{Medio} \\
R02 & Filtraci\'on de datos de facturaci\'on y obras sociales & 2 $\times$ 4 = 8 & \cellcolor{nivelmedio}\textbf{Medio} \\
\bottomrule
\end{tabularx}
\end{center}

{\small\color{teal!70!black}Clasificaci\'on seg\'un la matriz de la c\'atedra: cr\'itico 16--25, alto 10--15, medio 5--9, bajo 1--4.}

\newpage

\section*{Estado de los planes de acci\'on}
\addcontentsline{toc}{section}{Estado de los planes de accion}

Los 3 planes se encuentran en estado \textbf{Mitigation Planned}, con 0\% de avance.
El riesgo residual solo podr\'a reconocerse una vez implementados y verificados los
controles propuestos.

\begin{center}
\small
\renewcommand{\arraystretch}{1.6}
\begin{tabularx}{\textwidth}{@{} X >{\centering\arraybackslash}p{1.3cm} >{\raggedright\arraybackslash}p{3.4cm} >{\centering\arraybackslash}p{2.1cm} >{\centering\arraybackslash}p{1.8cm} @{}}
\toprule
\textbf{Plan} & \textbf{Riesgo} & \textbf{Responsable} & \textbf{Vencimiento} & \textbf{Presupuesto} \\
\midrule
Control de acceso por rol / especialidad & R01 & Jefe de Sistemas & 05/10/2026 & Bajo \\
Segmentaci\'on de red para equipos IoT & R06 & Jefe de Sistemas & 20/10/2026 & Medio \\
Pol\'itica de contrase\~nas obligatoria + MFA & R07 & Jefe de Sistemas / Admin.\ IT & 28/09/2026 & Medio \\
\bottomrule
\end{tabularx}
\end{center}

{\small\color{teal!70!black}Nota: R02 y R03 (nivel Medio) a\'un no cuentan con un plan formal asociado.}

\vspace{6mm}

\begin{decisionbox}[Seguimiento sugerido]
Solicitar una actualizaci\'on mensual del porcentaje de avance de cada plan al
responsable asignado, y evaluar en 60 d\'ias si corresponde definir planes formales
adicionales para R02 y R03.
\end{decisionbox}

\vspace{8mm}

\section*{Recomendaciones prioritarias}
\addcontentsline{toc}{section}{Recomendaciones prioritarias}

\begin{enumerate}[leftmargin=6mm,itemsep=5pt]
  \item \textbf{Implementar cuanto antes una pol\'itica de contrase\~nas obligatoria y MFA.}
  Es el riesgo de mayor nivel (R07, Cr\'itico) y hoy no cuenta con ning\'un control.

  \item \textbf{Ejecutar y dar seguimiento a los 3 planes ya definidos.}
  Control de acceso por rol, segmentaci\'on de red para equipos m\'edicos conectados,
  y pol\'itica de contrase\~nas. Medir avance mensualmente.

  \item \textbf{Capacitar al personal en detecci\'on de correos fraudulentos (phishing).}
  Hoy no existe ning\'un entrenamiento al respecto, y es la puerta de entrada m\'as
  com\'un para el robo de credenciales.

  \item \textbf{Revisar peri\'odicamente los permisos de acceso a facturaci\'on y obras sociales.}
  Para mantener vigente la separaci\'on de funciones ya implementada entre Analista y
  Auditor.

  \item \textbf{Presupuestar una copia de respaldo externa (fuera del sitio).}
  Para no depender de un \'unico punto de falla si el servidor principal sufre un
  incidente grave.
\end{enumerate}

\vfill
\begin{center}
{\scriptsize\color{teal!70!black}
Documento elaborado en el marco del Trabajo Pr\'actico de Seguridad de Sistemas (UCH).
Todos los datos personales y organizacionales utilizados son ficticios.}
\end{center}

\end{document}